%!TEX root = ../dissertation.tex
\chapter{Supplementary Analysis}
\label{AppendixB}

This appendix provides supplementary analyses referenced from the main text. These analyses support but are not required for the main arguments.

\section{Cross-Model Correlation Matrix}
\label{app:kendall-matrix}

Figure~\ref{fig:kendall-heatmap} reports the full $10 \times 10$ Kendall's $\tau$ rank-correlation matrix for per-prompt binary hallucination indicators on the 449-prompt initial benchmark. Each cell represents the ordinal association between two models' hallucination patterns: $\tau = 1$ would indicate that every prompt hallucinated by one model is also hallucinated by the other; $\tau = 0$ indicates no association.

\begin{figure}[ht]
\centering
\includegraphics[width=0.85\textwidth]{figures/consistency_heatmap.png}
\caption{Kendall's $\tau$ rank correlation of per-prompt binary hallucination indicators across all 10 models on the 449-prompt initial benchmark. Mean $\tau = 0.319$ (SD $= 0.137$, range 0.068--0.617). Of 45 model pairs, 42 are significant at $p < 0.05$; the three non-significant pairs all involve Qwen~3~Next~80B. Same-family pairs (e.g., GPT-4.1-mini / GPT-4o-mini, $\tau = 0.617$) show stronger correlations than cross-family pairs.}
\label{fig:kendall-heatmap}
\end{figure}

The mean pairwise correlation is $\tau = 0.319$ (median 0.303, SD $= 0.137$). Of the 45 unique pairs, 42 are statistically significant ($p < 0.05$); the three non-significant pairs all involve Qwen~3~Next~80B (vs.\ Claude Sonnet: $\tau = 0.068$, $p = 0.150$; vs.\ GPT-4.1: $\tau = 0.068$, $p = 0.149$; vs.\ Mixtral: $\tau = 0.076$, $p = 0.108$). The strongest correlations occur within the OpenAI family: GPT-4.1-mini and GPT-4o-mini ($\tau = 0.617$), GPT-4.1 and GPT-4.1-mini ($\tau = 0.524$), and GPT-4.1 and GPT-4o-mini ($\tau = 0.482$). One cross-family pair shows unexpectedly high agreement: Claude Sonnet~4.5 and GPT-5.1 ($\tau = 0.590$), suggesting shared training-data effects that cross provider boundaries.

The moderate mean correlation indicates that hallucination difficulty is partly a property of the prompt itself---not solely determined by model-specific factors. This shared difficulty signal motivates the geometric analysis in Chapter~\ref{ch:geo-predict-hallucination}: if prompts carry model-agnostic risk, that risk may be legible from the prompt's position in embedding space.

\section{Embedding Model Robustness}
\label{app:embedding-robustness}

All geometric analyses in this thesis use \texttt{text-embedding-3-large} (3,072 dimensions; OpenAI). To assess whether the hallucination-geometry relationship depends on this specific embedding model, we conducted a preliminary comparison using three embedding models on the 449-prompt initial benchmark:

\begin{itemize}[nosep]
    \item \texttt{text-embedding-3-small} (1,536 dimensions; OpenAI)---used in the initial study before upgrading to the large variant.
    \item \texttt{text-embedding-3-large} (3,072 dimensions; OpenAI)---the model used throughout the thesis.
    \item \texttt{all-mpnet-base-v2} (768 dimensions; open-source, Sentence Transformers).
\end{itemize}

For each embedding model, we computed the same geometric features (density, centrality, curvature) on 368 prompts from the four main categories and measured their Spearman rank correlation with binary hallucination indicators. Table~\ref{tab:embedding-robustness} reports the results; Figure~\ref{fig:embedding-robustness} visualizes the comparison.

\begin{table}[ht]
\centering
\small
\begin{tabular}{llrr}
\toprule
\textbf{Embedding Model} & \textbf{Feature} & \textbf{Spearman $r$} & \textbf{$p$-value} \\
\midrule
\texttt{text-embedding-3-small} & centrality & $-0.116$ & $1.9 \times 10^{-12}$ \\
 & density & $+0.091$ & $3.1 \times 10^{-8}$ \\
 & curvature & $-0.005$ & 0.776 \\
\midrule
\texttt{text-embedding-3-large} & centrality & $-0.046$ & 0.002 \\
 & density & $+0.011$ & 0.451 \\
 & curvature & $-0.035$ & 0.020 \\
\midrule
\texttt{all-mpnet-base-v2} & centrality & $-0.111$ & $7.1 \times 10^{-14}$ \\
 & density & $-0.003$ & 0.827 \\
 & curvature & $+0.006$ & 0.673 \\
\bottomrule
\end{tabular}
\caption{Spearman rank correlations between geometric features and binary hallucination indicators across three embedding models on the initial benchmark (368 prompts, four main categories). These are \emph{overall} correlations that conflate category structure; the within-category analysis in Chapter~\ref{ch:geo-predict-hallucination} uses the expanded benchmark with category-level controls.}
\label{tab:embedding-robustness}
\end{table}

\begin{figure}[ht]
\centering
\includegraphics[width=0.75\textwidth]{figures/embedding_robustness_comparison.png}
\caption{Spearman correlations between geometric features and hallucination across three embedding models. Centrality shows the most stable relationship: negative correlations across all three models, significant in two of three. Density and curvature are less stable, with sign and significance varying across models.}
\label{fig:embedding-robustness}
\end{figure}

Two patterns emerge. First, \textbf{centrality is the most robust feature}: it shows a negative correlation with hallucination (prompts farther from the centroid are more likely to hallucinate) across all three embedding models, reaching significance in two of three ($p < 10^{-12}$ for small and mpnet; $p = 0.002$ for large). Second, \textbf{density is embedding-model-dependent}: the significant positive correlation found with \texttt{text-embedding-3-small} ($r = +0.091$, $p < 10^{-7}$) does not replicate with the other two models. Curvature is near zero and non-significant across all three models.

These results carry two important caveats. First, this comparison uses overall (not within-category) correlations on the smaller initial benchmark, which conflates the category-structure signal that Chapter~\ref{ch:geo-predict-hallucination} carefully separates. The within-category density signal---the thesis's central geometric finding---was established on the expanded benchmark using \texttt{text-embedding-3-large} and has not been tested with alternative embedding models. Second, oppositeness was excluded from this comparison because it depends on PCA axes that are inherently embedding-model-specific, making cross-model comparison uninformative.

This preliminary check suggests that centrality captures a stable structural property of prompt embeddings, while density's hallucination signal is more sensitive to the embedding space in which it is computed. A full robustness analysis---recomputing within-category features on the expanded benchmark with multiple embedding models---remains a direction for future work (Section~\ref{sec:ch5-limitations}).
